% Options for packages loaded elsewhere
\PassOptionsToPackage{unicode}{hyperref}
\PassOptionsToPackage{hyphens}{url}
%
\documentclass[
]{article}
\usepackage{amsmath,amssymb}
\usepackage{iftex}
\ifPDFTeX
  \usepackage[T1]{fontenc}
  \usepackage[utf8]{inputenc}
  \usepackage{textcomp} % provide euro and other symbols
\else % if luatex or xetex
  \usepackage{unicode-math} % this also loads fontspec
  \defaultfontfeatures{Scale=MatchLowercase}
  \defaultfontfeatures[\rmfamily]{Ligatures=TeX,Scale=1}
\fi
\usepackage{lmodern}
\ifPDFTeX\else
  % xetex/luatex font selection
\fi
% Use upquote if available, for straight quotes in verbatim environments
\IfFileExists{upquote.sty}{\usepackage{upquote}}{}
\IfFileExists{microtype.sty}{% use microtype if available
  \usepackage[]{microtype}
  \UseMicrotypeSet[protrusion]{basicmath} % disable protrusion for tt fonts
}{}
\makeatletter
\@ifundefined{KOMAClassName}{% if non-KOMA class
  \IfFileExists{parskip.sty}{%
    \usepackage{parskip}
  }{% else
    \setlength{\parindent}{0pt}
    \setlength{\parskip}{6pt plus 2pt minus 1pt}}
}{% if KOMA class
  \KOMAoptions{parskip=half}}
\makeatother
\usepackage{xcolor}
\usepackage{longtable,booktabs,array}
\usepackage{calc} % for calculating minipage widths
% Correct order of tables after \paragraph or \subparagraph
\usepackage{etoolbox}
\makeatletter
\patchcmd\longtable{\par}{\if@noskipsec\mbox{}\fi\par}{}{}
\makeatother
% Allow footnotes in longtable head/foot
\IfFileExists{footnotehyper.sty}{\usepackage{footnotehyper}}{\usepackage{footnote}}
\makesavenoteenv{longtable}
\setlength{\emergencystretch}{3em} % prevent overfull lines
\providecommand{\tightlist}{%
  \setlength{\itemsep}{0pt}\setlength{\parskip}{0pt}}
\setcounter{secnumdepth}{-\maxdimen} % remove section numbering
\ifLuaTeX
\usepackage[bidi=basic]{babel}
\else
\usepackage[bidi=default]{babel}
\fi
\babelprovide[main,import]{french}
% get rid of language-specific shorthands (see #6817):
\let\LanguageShortHands\languageshorthands
\def\languageshorthands#1{}
\ifLuaTeX
  \usepackage{selnolig}  % disable illegal ligatures
\fi
\usepackage{bookmark}
\IfFileExists{xurl.sty}{\usepackage{xurl}}{} % add URL line breaks if available
\urlstyle{same}
\hypersetup{
  pdftitle={Étude de faisabilité technique et juridique},
  pdfauthor={Équipe GOVIA Cameroun},
  pdflang={fr},
  hidelinks,
  pdfcreator={LaTeX via pandoc}}

\title{Étude de faisabilité technique et juridique}
\usepackage{etoolbox}
\makeatletter
\providecommand{\subtitle}[1]{% add subtitle to \maketitle
  \apptocmd{\@title}{\par {\large #1 \par}}{}{}
}
\makeatother
\subtitle{GOVIA Cameroun --- Agent intelligent d'assistance
administrative et de conformité aux procédures publiques}
\author{Équipe GOVIA Cameroun}
\date{26 octobre 2026}

\begin{document}
\maketitle

\section{GOVIA Cameroun --- Étude de faisabilité technique et
juridique}\label{govia-cameroun-uxe9tude-de-faisabilituxe9-technique-et-juridique}

\textbf{Agent intelligent d'assistance administrative et de conformité
aux procédures publiques}

\begin{center}\rule{0.5\linewidth}{0.5pt}\end{center}

\subsection{1. Synthèse exécutive}\label{synthuxe8se-exuxe9cutive}

GOVIA Cameroun est un \textbf{agent d'IA agentique bilingue (français /
anglais)} qui assiste le citoyen dans ses démarches administratives : il
identifie la procédure, vérifie les pièces requises, explique les
étapes, \textbf{cite les textes officiels applicables} et propose de
\textbf{pré-remplir le formulaire officiel} pour que le citoyen le
valide et le soumette. Un \textbf{second agent (« auditeur de conformité
»)} contrôle que chaque réponse et action de l'assistant est justifiée
par un texte de la base de connaissances, et détecte les incohérences
entre procédures et textes.

\textbf{Conclusion de faisabilité : OUI, réalisable --- à condition de
respecter quatre contraintes majeures.}

\begin{longtable}[]{@{}
  >{\raggedright\arraybackslash}p{(\columnwidth - 4\tabcolsep) * \real{0.3333}}
  >{\raggedright\arraybackslash}p{(\columnwidth - 4\tabcolsep) * \real{0.3333}}
  >{\raggedright\arraybackslash}p{(\columnwidth - 4\tabcolsep) * \real{0.3333}}@{}}
\toprule\noalign{}
\begin{minipage}[b]{\linewidth}\raggedright
Axe
\end{minipage} & \begin{minipage}[b]{\linewidth}\raggedright
Verdict
\end{minipage} & \begin{minipage}[b]{\linewidth}\raggedright
Conditions / risques majeurs
\end{minipage} \\
\midrule\noalign{}
\endhead
\bottomrule\noalign{}
\endlastfoot
Faisabilité juridique & OK Faisable & Respect de la \textbf{Loi n°
2024/017} (protection des données), consentement explicite, \textbf{pas
de transfert transfrontalier non autorisé}, pas de traitement de données
biométriques, responsabilité gardée par le citoyen. \\
Faisabilité technique & OK Faisable & Part de zéro, pile
\textbf{gratuite / low-cost} (LLM local ou API à tarification dégressive
+ RAG local), OCR local, base de connaissances ouverte. \\
Conformité hackathon & OK Respectable & Deux parties obligatoires
remplies (technique + réglementaire), pipeline \textbf{Y.3172}
cartographié, base RAG ouverte, \textbf{analyse d'écarts} produite,
politiques \textbf{réelles et ouvertes}, données de cas d'usage
\textbf{synthétiques}. \\
Calendrier & ATTENTION Urgent & Délai \textbf{10 novembre 2026}. Plan en
3 phases ciblé sur un MVP solide + démo. \\
\end{longtable}

\textbf{Recommandations clés :} 1. Adopter le \textbf{modèle local
(Ollama + Mistral/Qwen)} pour la conformité data ; utiliser une
\textbf{API cloud (Groq/DeepSeek)} avec \textbf{anonymisation stricte}
uniquement pour la démo rapide. 2. \textbf{Ne pas inventer de
politiques.} Documenter chaque lacune (gap) --- c'est la partie la plus
valorisée du hackathon. 3. La \textbf{CNI n'est qu'un exemple} :
construire une base de connaissances \textbf{multi-démarches}
(passeport, état civil, nationalité, contrat de travail, NIU, carte
grise, etc.). 4. Faire du \textbf{mapping Y.3172 détaillé} et de la
\textbf{comparaison version A (sans gouvernance) vs version B
(gouvernée)} les piliers de la démo.

\begin{center}\rule{0.5\linewidth}{0.5pt}\end{center}

\subsection{2. Contexte, objectifs et contraintes du
hackathon}\label{contexte-objectifs-et-contraintes-du-hackathon}

\subsubsection{2.1 Le hackathon ITU « AI for Good Lab »
(UAE)}\label{le-hackathon-itu-ai-for-good-lab-uae}

Activité crowdsourcée ouverte à tous (étudiants, chercheurs,
professionnels, startups). Objectif : construire, par pays, une
\textbf{base de connaissances ouverte} et des \textbf{agents d'IA
(agentic AI)} pour étudier comment le comportement des agents change
selon les politiques appliquées.

\textbf{Deux parties obligatoires (les deux requises pour concourir,
}soumises avant le 10 novembre\textbf{) :} - \textbf{Partie 1 ---
Solution technique} : \textbf{le système fonctionnel} (agent + RAG +
auditeur + démo). Le système \textbf{cite les textes} dans ses réponses
et \textbf{vérifie} sa conformité en interne. C'est le produit. -
\textbf{Partie 2 --- Ancrage réglementaire} : \textbf{le dossier} qui
trace, étape par étape du pipeline Y.3172, \textbf{quel texte réel
gouverne l'action}, et l'\textbf{analyse d'écarts} (où les textes
manquent). C'est la justification documentaire.

\begin{quote}
\textbf{À noter :} Partie 1 \textbf{contient} des justifications (le
système cite les textes en pratique), Partie 2 \textbf{formalise} cette
justification pour les juges. La distinction est \textbf{fonction (P1)
vs dossier (P2)}, pas « avec/sans justification ».
\end{quote}

\textbf{Ce qui est évalué :} - \textbf{Mapping détaillé sur le pipeline
ITU-T Y.3172} (critère le plus valorisé). - \textbf{Analyse d'écarts
(gap analysis)} --- partie la plus valorisée. - \textbf{Base de
connaissances / RAG} (partie intégrante de la solution). - Mapping aux
\textbf{dimensions de l'AI Readiness Report}. - Bonus : comparer
\textbf{version A (sans gouvernance) vs version B (gouvernée)} ;
comparer une lacune du pays avec la réglementation d'un autre pays.

\textbf{Ce qui n'est PAS évalué :} - \textbf{L'exactitude du modèle
(accuracy)} n'est pas un critère. - \textbf{La valeur business /
commerciale} n'est pas un critère. Un pitch deck sans prototype
fonctionnel est écarté.

\subsubsection{2.2 Contraintes et règles de conformité (à respecter
scrupuleusement)}\label{contraintes-et-ruxe8gles-de-conformituxe9-uxe0-respecter-scrupuleusement}

\begin{longtable}[]{@{}
  >{\raggedright\arraybackslash}p{(\columnwidth - 2\tabcolsep) * \real{0.5000}}
  >{\raggedright\arraybackslash}p{(\columnwidth - 2\tabcolsep) * \real{0.5000}}@{}}
\toprule\noalign{}
\begin{minipage}[b]{\linewidth}\raggedright
Règle
\end{minipage} & \begin{minipage}[b]{\linewidth}\raggedright
Implication pour GOVIA
\end{minipage} \\
\midrule\noalign{}
\endhead
\bottomrule\noalign{}
\endlastfoot
Interdiction formelle d'utiliser des informations confidentielles (accès
restreint) & Seules les \textbf{sources publiées et ouvertes} sont
admises. Pas d'accès à des systèmes IDCAM/DGI internes. \\
\textbf{Ne pas inventer de politiques} & En l'absence de document
applicable, \textbf{documenter la lacune (gap)} --- valorisé, pas
pénalisé. \\
Données du cas d'usage peuvent être
générées/augmentées/simulées/synthétiques & Le scénario citoyen, les
pièces, les formulaires sont \textbf{synthétiques} pour la démo. \\
Les documents de politique doivent être \textbf{réels --- jamais
fictifs} & La base de connaissances cite des textes réels (décrets,
lois) avec références. \\
Modèles propriétaires / non open source autorisés & Seule la
\textbf{base de connaissances doit être ouverte} (le code peut rester
privé). \\
Base de connaissances : découpée en chunks machine-readable, embarquée
dans la solution (RAG) pour traçabilité & Chaque décision de l'agent
doit être \textbf{justifiable par une politique}. \\
Taille modeste attendue (\textasciitilde30 politiques ou moins) &
Suffisant si les politiques gouvernent réellement chaque étape du
pipeline. \\
Code source peut rester privé ; la base de connaissances doit rester
ouverte & Livraison : rapport (5 pages max), dépôt GitHub, base de
connaissances, vidéo démo (7 min max). \\
Composition d'équipe : max 4 (4 participants ou 3 + 1 mentor) & À
vérifier pour l'inscription. \\
Date critique : \textbf{10 novembre 2026} (soumission complète) & Plan
d'action pragmatique (§7). \\
\end{longtable}

\subsubsection{2.3 Pipeline ITU-T Y.3172 --- élément
central}\label{pipeline-itu-t-y.3172-uxe9luxe9ment-central}

Structure attendue : un \textbf{bloc intent} (le prompt) → un
\textbf{orchestrateur} → le pipeline lui-même, dupliqué en deux
instances identiques : un \textbf{sandbox} (test/validation) et
l'\textbf{environnement réseau de production}.

\begin{longtable}[]{@{}
  >{\raggedright\arraybackslash}p{(\columnwidth - 4\tabcolsep) * \real{0.3333}}
  >{\raggedright\arraybackslash}p{(\columnwidth - 4\tabcolsep) * \real{0.3333}}
  >{\raggedright\arraybackslash}p{(\columnwidth - 4\tabcolsep) * \real{0.3333}}@{}}
\toprule\noalign{}
\begin{minipage}[b]{\linewidth}\raggedright
Nœud
\end{minipage} & \begin{minipage}[b]{\linewidth}\raggedright
Signification
\end{minipage} & \begin{minipage}[b]{\linewidth}\raggedright
Ce qu'il faut décrire
\end{minipage} \\
\midrule\noalign{}
\endhead
\bottomrule\noalign{}
\endlastfoot
\textbf{SRC} & Source de données & Nature des données, métadonnées,
technologie de collecte, emplacement, fréquence/durée, configuration. \\
\textbf{C} & Collecteur & Comment les données sont agrégées : plateforme
edge, transmission réseau, cloud cœur ou edge. \\
\textbf{PP} & Pré-traitement & Nettoyage, gestion multimodal,
standardisation des formats, \textbf{dé-identification/anonymisation},
traitement des valeurs nulles. \\
\textbf{M} & Modèle & Type (régression, détection d'anomalie, vision,
audio, multimodal\ldots), hébergement, paramètres, fonction. \\
\textbf{P} & \textbf{Politique} & À chaque étape, la \textbf{politique
nationale applicable} : loi de protection des données, niveau
d'anonymisation requis, etc. \\
\textbf{D} & Distributeur & Comment la sortie du modèle est acheminée
vers sa cible, identification de la cible. \\
\textbf{SINK} & Cible & Ce qui exécute le résultat (application,
actionneur, plateforme d'analyse\ldots). \\
\end{longtable}

\begin{center}\rule{0.5\linewidth}{0.5pt}\end{center}

\subsection{3. Description du cas d'usage et concept de
solution}\label{description-du-cas-dusage-et-concept-de-solution}

\subsubsection{3.1 Problème}\label{probluxe8me}

Le citoyen ne maîtrise pas la procédure exacte d'une démarche
administrative (ex. renouvellement de CNI), ignore les \textbf{pièces
exigibles}, et les dossiers incomplets ou incorrects coûtent temps,
argent et refus. L'information administrative est \textbf{dispersée en
silos} (IDCAM, PassCam, DGI, SID, mairies\ldots) et peu accessible.

\subsubsection{3.2 Solution --- GÉNÉRALISTE (la CNI n'est qu'un
exemple)}\label{solution-guxe9nuxe9raliste-la-cni-nest-quun-exemple}

\begin{quote}
ATTENTION \textbf{Clarification importante :} la procédure CNI illustrée
dans le scénario initial est \textbf{un exemple parmi beaucoup
d'autres}. GOVIA couvre \textbf{le maximum de démarches administratives
possibles} (voir §6).
\end{quote}

\textbf{GOVIA} est un agent bilingue (FR/EN) qui, à partir de la
question du citoyen : 1. \textbf{Identifie} la procédure administrative
concernée. 2. \textbf{Vérifie les pièces requises} (selon la démarche).
3. \textbf{Explique les étapes} de la procédure. 4. \textbf{Cite les
textes officiels applicables} (décret, loi, arrêté) avec référence
précise. 5. Propose de \textbf{pré-remplir le formulaire officiel} ; le
citoyen \textbf{valide et soumet}.

\subsubsection{3.3 Architecture en deux
agents}\label{architecture-en-deux-agents}

\begin{longtable}[]{@{}
  >{\raggedright\arraybackslash}p{(\columnwidth - 2\tabcolsep) * \real{0.5000}}
  >{\raggedright\arraybackslash}p{(\columnwidth - 2\tabcolsep) * \real{0.5000}}@{}}
\toprule\noalign{}
\begin{minipage}[b]{\linewidth}\raggedright
Agent
\end{minipage} & \begin{minipage}[b]{\linewidth}\raggedright
Rôle
\end{minipage} \\
\midrule\noalign{}
\endhead
\bottomrule\noalign{}
\endlastfoot
\textbf{GOVIA-Assistant} & Dialogue avec le citoyen ; identifie la
démarche ; vérifie les pièces ; explique les étapes ; cite les textes ;
pré-remplit le formulaire. \\
\textbf{GOVIA-Auditeur (conformité)} & Contrôle que chaque réponse et
action de l'assistant est \textbf{justifiée par un texte de la base de
connaissances} ; détecte les incohérences ; bloque/signale si une action
n'est pas couverte par un texte. \\
\end{longtable}

Cette architecture est directement inspirée du thème : « Un second agent
peut contrôler si les réponses respectent les règles publiées. »

\subsubsection{\texorpdfstring{3.4 Autonomie choisie :
\textbf{pré-remplissage} (citoyen valide et
soumet)}{3.4 Autonomie choisie : pré-remplissage (citoyen valide et soumet)}}\label{autonomie-choisie-pruxe9-remplissage-citoyen-valide-et-soumet}

L'agent \textbf{pré-remplit} le formulaire mais \textbf{ne soumet pas} à
la place du citoyen. Le citoyen garde la maîtrise et la responsabilité
de la soumission. Cela limite fortement le risque juridique (§4) et
respecte la règle du hackathon (le citoyen valide).

\begin{center}\rule{0.5\linewidth}{0.5pt}\end{center}

\subsection{4. Faisabilité juridique}\label{faisabilituxe9-juridique}

\subsubsection{4.1 Socle juridique camerounais --- réel et
ouvert}\label{socle-juridique-camerounais-ruxe9el-et-ouvert}

Tous les textes ci-dessous sont \textbf{réels, publics et ouverts} ---
ils constituent le cœur de la base de connaissances (Partie 2 du
hackathon).

\begin{longtable}[]{@{}
  >{\raggedright\arraybackslash}p{(\columnwidth - 6\tabcolsep) * \real{0.2500}}
  >{\raggedright\arraybackslash}p{(\columnwidth - 6\tabcolsep) * \real{0.2500}}
  >{\raggedright\arraybackslash}p{(\columnwidth - 6\tabcolsep) * \real{0.2500}}
  >{\raggedright\arraybackslash}p{(\columnwidth - 6\tabcolsep) * \real{0.2500}}@{}}
\toprule\noalign{}
\begin{minipage}[b]{\linewidth}\raggedright
Instrument
\end{minipage} & \begin{minipage}[b]{\linewidth}\raggedright
Référence
\end{minipage} & \begin{minipage}[b]{\linewidth}\raggedright
Pertinence pour GOVIA
\end{minipage} & \begin{minipage}[b]{\linewidth}\raggedright
Statut
\end{minipage} \\
\midrule\noalign{}
\endhead
\bottomrule\noalign{}
\endlastfoot
Décret n° 2016/375 du 04 août 2016 & CNI --- caractéristiques \&
modalités de délivrance (art. 5 obligation ; art. 11 pièces) & Texte
principal du scénario de référence CNI & OK Ouvert \\
Décret n° 2021/347 du 17 juin 2021 & Passeports biométriques ---
conditions d'établissement & Démarche passeport & OK Ouvert \\
Ordonnance n° 81/02 du 29 juin 1981 & Organisation de l'état civil
(naissance, mariage, décès) & Actes d'état civil & OK Ouvert \\
Loi n° 2011/011 du 06 mai 2011 & Délais de déclaration des actes d'état
civil & Délais naissance/décès & OK Ouvert \\
Code civil camerounais & État des personnes (art. 34, 49, 76, 78, 79) &
Actes d'état civil & OK Ouvert \\
Loi n° 1968-LF-3 du 11 juin 1968 & Code de la nationalité camerounaise
(art. 42 certificat) & Certificat de nationalité & OK Ouvert \\
Décret n° 68-DF-478 du 16 déc. 1968 & Modalités d'application du Code de
la nationalité & Certificat de nationalité & OK Ouvert \\
Loi n° 92-007 du 14 août 1992 & Code du travail (art. 23, 25 contrat
CDI/CDD) & Contrat de travail & OK Ouvert \\
\textbf{Loi n° 2024/017 du 23 déc. 2024} & \textbf{Protection des
données à caractère personnel} & Consentement, DPIA, \textbf{transferts
transfrontaliers (art. 32)}, données sensibles (art. 33), sanctions & OK
Ouvert \\
Loi n° 2010/012 du 21 déc. 2010 & Cybersécurité et cybercriminalité
(signature électronique, ANTIC) & Signature électronique, sécurité & OK
Ouvert \\
Loi n° 2010/021 du 21 déc. 2010 & Commerce électronique &
Documents/signatures électroniques & OK Ouvert \\
Loi n° 2010/013 du 21 déc. 2010 & Communications électroniques (modifiée
par 2015/006) & Contexte télécoms & OK Ouvert \\
OHADA --- Acte Uniforme AUDCG & Signature électronique (effet
supranational) & Validité signature électronique & OK Ouvert \\
Décret n° 79/341 du 03 sept. 1979 & Immatriculation des véhicules (carte
grise) & Carte grise & OK Ouvert \\
Décret n° 2020/PM du 01 juin 2020 & Documents de transport sécurisés
(cartes grises) & Carte grise moderne & OK Ouvert \\
Décret n° 2016/373 du 04 août 2016 & Entrée, séjour et sortie des
étrangers (loi 97/012) & Carte de séjour / étrangers & OK Ouvert \\
Procédure IDCAM (DGSN) & idcam.cm, 17 février 2025 --- pré-enrôlement
CNI & Procédure CNI réelle & OK Ouvert \\
Procédure PassCam (DGSN) & passcam.cm, 2024 --- pré-enrôlement passeport
& Procédure passeport réelle & OK Ouvert \\
FGI 2026 / E-Gov'A 2026 & DPI, interopérabilité, identité numérique &
Contexte stratégique & OK Ouvert \\
\end{longtable}

\subsubsection{4.2 Analyse par enjeu}\label{analyse-par-enjeu}

\paragraph{\texorpdfstring{a) Protection des données personnelles ---
\textbf{Loi n° 2024/017 du 23 décembre
2024}}{a) Protection des données personnelles --- Loi n° 2024/017 du 23 décembre 2024}}\label{a-protection-des-donnuxe9es-personnelles-loi-n-2024017-du-23-duxe9cembre-2024}

La Loi n° 2024/017 est \textbf{l'instrument le plus contraignant} pour
GOVIA. Points clés :

\begin{itemize}
\tightlist
\item
  \textbf{Consentement} : tout traitement de données personnelles
  requiert un consentement \textbf{clair, distinct, explicite} du
  citoyen. GOVIA doit recueillir ce consentement avant de traiter les
  pièces et avant le pré-remplissage.
\item
  \textbf{Minimisation et finalité} : ne traiter que les données
  nécessaires à la démarche demandée.
\item
  \textbf{Données sensibles (art. 33)} : les données biométriques
  (empreintes, photo d'identité biométrique) sont \textbf{interdites de
  traitement sauf autorisation}. \textbf{GOVIA ne doit pas traiter les
  données biométriques} --- uniquement les pièces justificatives
  administratives (acte de naissance, certificat de nationalité, etc.).
  La capture biométrique reste exclusivement du ressort de
  l'administration (DGSN).
\item
  \textbf{DPIA} : un traitement à haut risque (massification de données
  administratives) requiert une \textbf{analyse d'impact relative à la
  protection des données}.
\item
  \textbf{Transferts transfrontaliers (art. 32)} : tout transfert de
  données personnelles \textbf{hors du Cameroun} requiert une
  \textbf{autorisation préalable de l'Autorité de protection des
  données}. Un LLM cloud étranger (OpenAI, Groq, Mistral USA/EU)
  implique un transfert → soit autorisation, soit \textbf{traitement
  local}, soit \textbf{anonymisation/désidentification avant envoi}.
\item
  \textbf{Responsabilité conjointe} : le donneur et le destinataire sont
  \textbf{solidairement responsables} en cas de transfert illicite.
\item
  \textbf{Sanctions} : amendes de 1 à 1 000 millions FCFA selon le
  manquement ; peines d'emprisonnement jusqu'à 10 ans pour transferts
  illégaux. Période transitoire \textbf{expirée le 23 juin 2026} → la
  loi est pleinement applicable.
\end{itemize}

\paragraph{b) Identité, état civil et titres --- décrets et
ordonnances}\label{b-identituxe9-uxe9tat-civil-et-titres-duxe9crets-et-ordonnances}

\begin{itemize}
\tightlist
\item
  \textbf{CNI (décret 2016/375)} : la possession est obligatoire dès 18
  ans ; les pièces exigent un certificat de nationalité, un acte de
  naissance, etc. (art. 11). Le pré-enrôlement IDCAM est la procédure
  réelle. GOVIA cite ces exigences textuelles.
\item
  \textbf{État civil (ordonnance 81/02, loi 2011/011)} : actes de
  naissance, mariage, décès ; délais de déclaration (naissance 60 jours,
  décès 90 jours). Jugement supplétif pour naissance non déclarée à
  temps (art. 33).
\item
  \textbf{Nationalité (loi 1968-LF-3)} : certificat de nationalité
  délivré par le tribunal (art. 42), pièce clé de la première CNI.
\end{itemize}

\paragraph{c) Transactions électroniques et signature --- validité du
pré-remplissage}\label{c-transactions-uxe9lectroniques-et-signature-validituxe9-du-pruxe9-remplissage}

\begin{itemize}
\tightlist
\item
  \textbf{Loi n° 2010/012 (cybersécurité)} et \textbf{Loi n° 2010/021
  (commerce électronique)} : la \textbf{signature électronique avancée}
  a \textbf{la même valeur juridique} que la signature manuscrite (art.
  17). \textbf{OHADA AUDCG} : la signature électronique ne peut être
  refusée au seul motif de sa forme électronique.
\item
  \textbf{Implication pour GOVIA} : le pré-remplissage du formulaire par
  l'agent + \textbf{validation/signature électronique du citoyen} est
  juridiquement solide. L'agent ne signe pas ; le citoyen valide et
  signe. Cela sécurise la conformité.
\end{itemize}

\paragraph{d) Cybersécurité et
responsabilité}\label{d-cybersuxe9curituxe9-et-responsabilituxe9}

\begin{itemize}
\tightlist
\item
  L'agent est un \textbf{intermédiaire technique non délégué} de
  l'administration. En cas d'erreur de pré-remplissage, \textbf{le
  citoyen reste responsable} des informations qu'il valide et soumet.
  GOVIA doit afficher cette limitation clairement.
\item
  La \textbf{responsabilité de l'agent IA tiers} dans les démarches
  administratives camerounaises n'a \textbf{pas de cadre explicite} →
  \textbf{lacune réglementaire à documenter} (voir §4.3).
\end{itemize}

\paragraph{e) Interopérabilité et DPI --- contexte
stratégique}\label{e-interopuxe9rabilituxe9-et-dpi-contexte-stratuxe9gique}

\begin{itemize}
\tightlist
\item
  Le Cameroun promeut les \textbf{Infrastructures Publiques Numériques
  (DPI)} : identité numérique, paiements, échanges sécurisés de données,
  registres de référence (FGI 2026, E-Gov'A 2026). Le ministère de la
  Fonction publique (AIGLES) recense 400+ processus, 9 250 règles de
  gestion.
\item
  \textbf{Implication} : GOVIA s'inscrit dans la volonté
  d'interopérabilité, mais les systèmes restent en silos (IDCAM,
  PassCam, DGI, SID séparés) → \textbf{gap} à documenter.
\end{itemize}

\subsubsection{\texorpdfstring{4.3 Analyse d'écarts juridiques (GAP
ANALYSIS) --- \textbf{partie la plus
valorisée}}{4.3 Analyse d'écarts juridiques (GAP ANALYSIS) --- partie la plus valorisée}}\label{analyse-duxe9carts-juridiques-gap-analysis-partie-la-plus-valorisuxe9e}

\begin{longtable}[]{@{}
  >{\raggedright\arraybackslash}p{(\columnwidth - 6\tabcolsep) * \real{0.2500}}
  >{\raggedright\arraybackslash}p{(\columnwidth - 6\tabcolsep) * \real{0.2500}}
  >{\raggedright\arraybackslash}p{(\columnwidth - 6\tabcolsep) * \real{0.2500}}
  >{\raggedright\arraybackslash}p{(\columnwidth - 6\tabcolsep) * \real{0.2500}}@{}}
\toprule\noalign{}
\begin{minipage}[b]{\linewidth}\raggedright
\#
\end{minipage} & \begin{minipage}[b]{\linewidth}\raggedright
Écart (gap) identifié
\end{minipage} & \begin{minipage}[b]{\linewidth}\raggedright
Conséquence
\end{minipage} & \begin{minipage}[b]{\linewidth}\raggedright
Mitigation / recommandation
\end{minipage} \\
\midrule\noalign{}
\endhead
\bottomrule\noalign{}
\endlastfoot
1 & \textbf{Pas d'API publique officielle} pour IDCAM / PassCam &
L'agent \textbf{ne peut pas soumettre} directement ; le pré-remplissage
est une \textbf{reconstitution/simulation}. & Documenter le gap ;
positionner GOVIA comme assistant de pré-remplissage + validation
citoyenne. Ne pas simuler une soumission réelle. \\
2 & \textbf{Absence de cadre explicite} pour la responsabilité d'un
agent IA tiers dans les démarches administratives & En cas d'erreur,
répartition de la responsabilité incertaine. & Afficher que le citoyen
reste responsable ; ajouter une clause de non-responsabilité ; proposer
un « carnet d'audit » traçable. \\
3 & \textbf{Transfert transfrontalier des données} (art. 32, Loi
2024/017) & Un LLM cloud étranger transfère des données hors du Cameroun
→ \textbf{illégal sans autorisation}. & Utiliser un \textbf{modèle
local} (Ollama) pour la conformité ; ou anonymiser/désidentifier avant
envoi cloud ; ou obtenir l'autorisation de l'Autorité. \\
4 & \textbf{Données biométriques} (art. 33) & Empreintes / photo
biométrique = données sensibles interdites de traitement sauf
autorisation. & \textbf{Exclure toute capture/traitement biométrique} de
GOVIA. Laisser la capture à l'administration (DGSN). \\
5 & \textbf{Absence de texte d'application} de la Loi 2024/017 (décret /
lignes directrices de l'Autorité) & Applicabilité partielle ;
incertitude sur les DPIA et autorisations. & Documenter le gap ; se
conformer aux principes généraux (consentement, minimisation,
sécurité). \\
6 & \textbf{Absence de cadre spécifique aux agents IA / IA générative
dans les services publics} & Pas de règles de classement, de
surveillance, d'audit des agents. & Documenter le gap (forte valeur).
Comparer avec l'UE (AI Act) et l'Afrique du Sud (POPIA) / Maroc (loi
09-08) --- \textbf{bonus}. \\
7 & \textbf{Fragmentation / silos administratifs} (IDCAM, PassCam, DGI,
SID séparés) & Pas d'échange de données structuré ; l'agent doit
interroger plusieurs sources. & Documenter le gap (aligné sur FGI 2026 «
de la fragmentation à l'interopérabilité ») ; proposer une base de
connaissances unifiée. \\
8 & \textbf{Consentement} explicite requis & L'agent doit recueillir un
consentement clair et distinct. & Implémenter un consentement dédié
avant traitement et pré-remplissage. \\
\end{longtable}

\textbf{Comparaison avec un autre pays (bonus) :} le Cameroun n'a
\textbf{pas de cadre équivalent à l'AI Act de l'UE} (règles par
catégorie de risque, obligations pour systèmes à haut risque,
surveillance humaine) ni de loi sur l'IA générative. L'\textbf{Afrique
du Sud (POPIA)} et le \textbf{Maroc (loi 09-08, CNDP)} ont des régimes
de transfert transfrontalier plus développés. \textbf{Gap de
normalisation} à soumettre à l'ITU.

\subsubsection{4.4 Matrice risque /
conformité}\label{matrice-risque-conformituxe9}

\begin{longtable}[]{@{}
  >{\raggedright\arraybackslash}p{(\columnwidth - 4\tabcolsep) * \real{0.3333}}
  >{\raggedright\arraybackslash}p{(\columnwidth - 4\tabcolsep) * \real{0.3333}}
  >{\raggedright\arraybackslash}p{(\columnwidth - 4\tabcolsep) * \real{0.3333}}@{}}
\toprule\noalign{}
\begin{minipage}[b]{\linewidth}\raggedright
Risque
\end{minipage} & \begin{minipage}[b]{\linewidth}\raggedright
Niveau
\end{minipage} & \begin{minipage}[b]{\linewidth}\raggedright
Mesure de mitigation
\end{minipage} \\
\midrule\noalign{}
\endhead
\bottomrule\noalign{}
\endlastfoot
Transfert transfrontalier illégal des données & \textbf{Élevé} & Modèle
local (Ollama) ; ou anonymisation avant envoi cloud ; ou autorisation de
l'Autorité. \\
Traitement de données biométriques & \textbf{Élevé} & Exclure toute
donnée biométrique de GOVIA. \\
Consentement absent & \textbf{Moyen} & Implémenter un consentement
explicite avant traitement/pré-remplissage. \\
Hallucination / conseil non fondé & \textbf{Moyen} & RAG obligatoire ;
l'auditeur bloque si pas de texte ; le citoyen valide. \\
Responsabilité en cas d'erreur & \textbf{Moyen} & Citoyen responsable ;
clause de non-responsabilité ; journal d'audit. \\
Utilisation de données confidentielles & \textbf{Élevé} & Seules des
sources ouvertes et réelles ; données de cas d'usage synthétiques. \\
\end{longtable}

\begin{center}\rule{0.5\linewidth}{0.5pt}\end{center}

\subsection{5. Faisabilité technique}\label{faisabilituxe9-technique}

\subsubsection{5.1 Architecture
proposée}\label{architecture-proposuxe9e}

\begin{verbatim}
[Citoyen] --prompt bilingue--> [Intent / orchestrateur]
                                     |
            +------------------------+------------------------+
            |                                                |
      [GOVIA-Assistant]                              [GOVIA-Auditeur]
         |  RAG sur base politiques                    | vérifie chaque
         |  extraction intent + requêtage              | réponse/action vs textes
         |  génération réponse + pré-remplissage        | blocage si gap
         |  journal d'audit                              |
      [Formulaire pré-rempli] --validation citoyenne--> [SINK : soumission citoyenne]
\end{verbatim}

\subsubsection{5.2 Mapping ITU-T Y.3172 détaillé
(esquisse)}\label{mapping-itu-t-y.3172-duxe9tailluxe9-esquisse}

\begin{longtable}[]{@{}
  >{\raggedright\arraybackslash}p{(\columnwidth - 4\tabcolsep) * \real{0.3333}}
  >{\raggedright\arraybackslash}p{(\columnwidth - 4\tabcolsep) * \real{0.3333}}
  >{\raggedright\arraybackslash}p{(\columnwidth - 4\tabcolsep) * \real{0.3333}}@{}}
\toprule\noalign{}
\begin{minipage}[b]{\linewidth}\raggedright
Nœud
\end{minipage} & \begin{minipage}[b]{\linewidth}\raggedright
Rôle dans GOVIA
\end{minipage} & \begin{minipage}[b]{\linewidth}\raggedright
Description détaillée
\end{minipage} \\
\midrule\noalign{}
\endhead
\bottomrule\noalign{}
\endlastfoot
\textbf{SRC} & Source de données & 1) Prompt du citoyen (texte, FR/EN) ;
2) pièces uploadées (images de documents : acte de naissance,
certificat, CNI\ldots) ; 3) \textbf{base de politiques ouvertes et
réelles} (décrets, lois, procédures). Métadonnées : source, date de
publication, autorité, URL. \\
\textbf{C} & Collecteur & Collecte de la conversation ; OCR des pièces
(Tesseract local ou vision LLM cloud avec anonymisation) ; récupération
des politiques pertinentes via indexation. \\
\textbf{PP} & Pré-traitement & Nettoyage du texte ; \textbf{détection de
langue (FR/EN)} ; \textbf{anonymisation / désidentification} des données
personnelles avant ingestion dans le modèle (masquage des noms, dates,
numéros) ; normalisation des formats (PDF, image, Word). \\
\textbf{M} & Modèle & LLM (multimodal si OCR de pièces) + \textbf{RAG}
sur la base politiques ; extraction d'intention (identifier la démarche)
+ requêtage des textes applicables ; génération de réponse +
pré-remplissage du formulaire. Hébergement : local (Ollama) ou cloud
(Groq/DeepSeek) --- voir §5.3. \\
\textbf{P} & \textbf{Politique} & À chaque étape, la politique
applicable : \textbf{Loi 2024/017} (protection des données,
anonymisation, consentement) ; \textbf{décret 2016/375} (CNI, pièces) ;
\textbf{ordonnance 81/02} (état civil) ; \textbf{loi 92-007} (contrat) ;
etc. L'agent cite la référence. \\
\textbf{D} & Distributeur & La sortie (réponse + formulaire pré-rempli +
rapport de conformité) est acheminée au citoyen via l'interface ; le
formulaire est présenté pour validation. \\
\textbf{SINK} & Cible & Le \textbf{citoyen} qui valide et soumet le
formulaire ; et l'\textbf{auditeur de conformité} qui reçoit le rapport
d'audit. \\
\end{longtable}

\subsubsection{5.3 Modèles : analyse cloud vs local (comme
demandé)}\label{moduxe8les-analyse-cloud-vs-local-comme-demanduxe9}

\begin{longtable}[]{@{}
  >{\raggedright\arraybackslash}p{(\columnwidth - 4\tabcolsep) * \real{0.3333}}
  >{\raggedright\arraybackslash}p{(\columnwidth - 4\tabcolsep) * \real{0.3333}}
  >{\raggedright\arraybackslash}p{(\columnwidth - 4\tabcolsep) * \real{0.3333}}@{}}
\toprule\noalign{}
\begin{minipage}[b]{\linewidth}\raggedright
Critère
\end{minipage} & \begin{minipage}[b]{\linewidth}\raggedright
\textbf{Cloud} (Groq, DeepSeek, OpenAI, Mistral API)
\end{minipage} & \begin{minipage}[b]{\linewidth}\raggedright
\textbf{Local} (Ollama + Mistral/Qwen/Llama)
\end{minipage} \\
\midrule\noalign{}
\endhead
\bottomrule\noalign{}
\endlastfoot
\textbf{Coût} & OK Tarification dégressive ; Groq/DeepSeek très
abordables ; tier gratuit pour démo & OK Gratuit après téléchargement
des modèles \\
\textbf{Conformité (Loi 2024/017 art. 32)} & NON Transfert
transfrontalier des données → autorisation requise ou anonymisation
stricte & OK \textbf{Aucune donnée ne sort du Cameroun} --- conformité
native \\
\textbf{Qualité linguistique (FR/EN)} & OK Très bonne ; modèles
multilingues performants & OK Bonne ; Mistral 7B / Qwen2.5 multilingues
acceptables \\
\textbf{OCR / multimodal} & OK Vision LLM cloud performant & ATTENTION
Possible (modèle de vision local) ou Tesseract local \\
\textbf{Matériel} & OK Aucun GPU local nécessaire & ATTENTION Besoin de
8--16 GB RAM ; fonctionne sur CPU \\
\textbf{Souveraineté / autonomie} & NON Dépendance réseau + fournisseur
& OK Souveraineté des données \\
\textbf{Latence} & OK Rapide & ATTENTION Légèrement plus lent (sans
GPU) \\
\textbf{Recommandation} & Démo rapide \textbf{avec anonymisation
stricte} & Version conformité recommandée \\
\end{longtable}

\textbf{Recommandation :} - \textbf{Prototype / démo rapide} :
\textbf{cloud (Groq ou DeepSeek API)} avec
\textbf{anonymisation/désidentification obligatoire avant envoi} et
\textbf{pas de données personnelles brutes}. Cela permet un déploiement
rapide sans matériel. - \textbf{Version conforme (production)} :
\textbf{modèle local (Ollama + Mistral 7B / Qwen2.5)}. Aucune donnée ne
sort du Cameroun → respect de l'art. 32 de la Loi 2024/017. Pour l'OCR,
utiliser \textbf{Tesseract local} ou un modèle de vision local.

\subsubsection{5.4 Base de connaissances --- RAG / graphe,
machine-readable,
ouverte}\label{base-de-connaissances-rag-graphe-machine-readable-ouverte}

\begin{itemize}
\tightlist
\item
  \textbf{Format} : chunks \textbf{machine-readable} (JSON/JSONL)
  découpés \textbf{par article/section} de chaque texte.
\item
  \textbf{Structure} : chaque chunk =
  \texttt{\{id,\ texte\_civil,\ texte\_officiel,\ source,\ URL,\ autorité,\ date,\ procedure,\ etape,\ tags\}}.
\item
  \textbf{Moteur} : \textbf{RAG vectoriel} (ChromaDB local + embeddings)
  + requêtage par similarité ; option \textbf{graphe de connaissances}
  (politique → démarche → étape → pièce) pour la traçabilité.
\item
  \textbf{Traçabilité} : chaque réponse de l'assistant doit citer les
  chunks utilisés → \textbf{journal d'audit}.
\item
  \textbf{Ouverture} : la base est \textbf{ouverte} (GitHub),
  versionnée, extensible à d'autres démarches.
\item
  \textbf{Taille} : \textasciitilde20--30 politiques suffisent pour
  couvrir les démarches du catalogue (§6).
\end{itemize}

\subsubsection{5.5 Comparaison version A (sans gouvernance) vs version B
(gouvernée) --- bonus
majeur}\label{comparaison-version-a-sans-gouvernance-vs-version-b-gouvernuxe9e-bonus-majeur}

\begin{longtable}[]{@{}
  >{\raggedright\arraybackslash}p{(\columnwidth - 4\tabcolsep) * \real{0.3333}}
  >{\raggedright\arraybackslash}p{(\columnwidth - 4\tabcolsep) * \real{0.3333}}
  >{\raggedright\arraybackslash}p{(\columnwidth - 4\tabcolsep) * \real{0.3333}}@{}}
\toprule\noalign{}
\begin{minipage}[b]{\linewidth}\raggedright
\end{minipage} & \begin{minipage}[b]{\linewidth}\raggedright
Version A (sans gouvernance)
\end{minipage} & \begin{minipage}[b]{\linewidth}\raggedright
Version B (gouvernée par les politiques)
\end{minipage} \\
\midrule\noalign{}
\endhead
\bottomrule\noalign{}
\endlastfoot
Requêtage de la base & Non systématique & \textbf{Obligatoire} ; chaque
réponse citée \\
Traçabilité & Aucune & \textbf{Journal d'audit} + références aux
textes \\
Hallucination & Risque élevé ; réponse non fondée & Bloquée par
l'auditeur si pas de texte \\
Incohérences & Non détectées & \textbf{Détectées et signalées} par
l'auditeur \\
Responsabilité & Incertaine & Citoyen responsable + consentement +
audit \\
Démo & Réponse rapide mais potentiellement erronée & Réponse fondée,
citable, vérifiée \\
\end{longtable}

\textbf{Scénario de démo :} poser une question ambiguë (ex. « Je veux
renouveler ma carte » sans préciser) → Version A donne une réponse
directe (potentiellement fausse) ; Version B demande clarification +
identifie les démarches possibles (CNI, carte grise, carte de séjour) +
cite les textes + l'auditeur signale l'ambiguïté.

\subsubsection{5.6 Stack recommandée (gratuite / low-cost, partir de
zéro)}\label{stack-recommanduxe9e-gratuite-low-cost-partir-de-zuxe9ro}

\begin{longtable}[]{@{}
  >{\raggedright\arraybackslash}p{(\columnwidth - 2\tabcolsep) * \real{0.5000}}
  >{\raggedright\arraybackslash}p{(\columnwidth - 2\tabcolsep) * \real{0.5000}}@{}}
\toprule\noalign{}
\begin{minipage}[b]{\linewidth}\raggedright
Composant
\end{minipage} & \begin{minipage}[b]{\linewidth}\raggedright
Recommandation
\end{minipage} \\
\midrule\noalign{}
\endhead
\bottomrule\noalign{}
\endlastfoot
Langage & Python 3.10+ \\
LLM & \textbf{Ollama + Mistral 7B / Qwen2.5} (conformité) ; ou
\textbf{Groq/DeepSeek API} (démo avec anonymisation) \\
RAG & \textbf{ChromaDB} (vectoriel local) + embeddings (HuggingFace
all-MiniLM ou local) ; option \textbf{FAISS} \\
OCR & \textbf{Tesseract} (local) ou vision LLM cloud (avec
anonymisation) \\
Chunking & Découpage par article/section → JSON/JSONL \\
Orchestration & \textbf{LangChain / LlamaIndex} (ou simplicité : module
custom) \\
API / interface & \textbf{FastAPI} (backend) + \textbf{Next.js / React +
Tailwind + next-intl} (frontend FR/EN) \\
Base de connaissances & JSON/JSONL + RAG, \textbf{ouverte}, versionnée
sur GitHub \\
Conteneurisation & Docker \\
Suivi / audit & Journal d'audit en JSON \\
\end{longtable}

\subsubsection{5.7 Faisabilité et risques
techniques}\label{faisabilituxe9-et-risques-techniques}

\begin{itemize}
\tightlist
\item
  \textbf{Partir de zéro} : le MVP est réalisable en \textasciitilde1
  semaine avec une pile simple (LLM + ChromaDB + Next.js + FastAPI).
\item
  \textbf{Multilinguisme FR/EN} : les modèles modernes (Mistral, Qwen,
  Groq) gèrent bien le mix FR/EN ; test requis.
\item
  \textbf{OCR des pièces} : Tesseract local gère les documents imprimés
  ; prévoir un étalonnage sur des pièces camerounaises.
\item
  \textbf{Pré-remplissage du formulaire} : reconstitution du formulaire
  IDCAM/PassCam \textbf{sans API réelle} → simulation de pré-remplissage
  (gap documenté).
\item
  \textbf{Scalabilité} : la base RAG est extensible à d'autres démarches
  (dimension cross-domain de l'AI Readiness).
\end{itemize}

\begin{center}\rule{0.5\linewidth}{0.5pt}\end{center}

\subsection{6. Catalogue des démarches administratives camerounaises
(section
dédiée)}\label{catalogue-des-duxe9marches-administratives-camerounaises-section-duxe9diuxe9e}

Recherche des démarches + leurs \textbf{textes réels et ouverts}.
Statuts : OK = texte ouvert identifié ; ATTENTION = texte partiellement
ouvert / gap.

\begin{longtable}[]{@{}
  >{\raggedright\arraybackslash}p{(\columnwidth - 10\tabcolsep) * \real{0.1667}}
  >{\raggedright\arraybackslash}p{(\columnwidth - 10\tabcolsep) * \real{0.1667}}
  >{\raggedright\arraybackslash}p{(\columnwidth - 10\tabcolsep) * \real{0.1667}}
  >{\raggedright\arraybackslash}p{(\columnwidth - 10\tabcolsep) * \real{0.1667}}
  >{\raggedright\arraybackslash}p{(\columnwidth - 10\tabcolsep) * \real{0.1667}}
  >{\raggedright\arraybackslash}p{(\columnwidth - 10\tabcolsep) * \real{0.1667}}@{}}
\toprule\noalign{}
\begin{minipage}[b]{\linewidth}\raggedright
\#
\end{minipage} & \begin{minipage}[b]{\linewidth}\raggedright
Démarche
\end{minipage} & \begin{minipage}[b]{\linewidth}\raggedright
Autorité compétente
\end{minipage} & \begin{minipage}[b]{\linewidth}\raggedright
Texte réel (référence)
\end{minipage} & \begin{minipage}[b]{\linewidth}\raggedright
Pièces clés
\end{minipage} & \begin{minipage}[b]{\linewidth}\raggedright
Statut
\end{minipage} \\
\midrule\noalign{}
\endhead
\bottomrule\noalign{}
\endlastfoot
1 & \textbf{CNI (renouvellement)} --- scénario de référence & DGSN
(idcam.cm) & Décret n° 2016/375 du 04 août 2016 (art. 5, 11) & Ancienne
CNI, acte de naissance, certificat de nationalité, acte de mariage
(femmes), justificatif de profession & OK \\
2 & \textbf{Passeport biométrique} & DGSN (passcam.cm) & Décret n°
2021/347 du 17 juin 2021 & Acte de naissance, certificat de nationalité,
CNI, jugement supplétif & OK \\
3 & \textbf{Acte de naissance} & Officier d'état civil (mairie) &
Ordonnance n° 81/02 du 29 juin 1981 ; Loi n° 2011/011 du 06 mai 2011 &
Déclaration ; jugement supplétif si tardif (art. 33) & OK \\
4 & \textbf{Acte de mariage} & Officier d'état civil & Ordonnance n°
81/02 (art. 48, 49) ; Code civil (art. 76) & Acte de naissance,
consentement, témoins & OK \\
5 & \textbf{Acte de décès} & Officier d'état civil & Ordonnance n° 81/02
(art. 78, 79) & Déclaration dans les 90 jours & OK \\
6 & \textbf{Certificat de nationalité} & Tribunal de première instance &
Loi n° 1968-LF-3 du 11 juin 1968 (art. 42) ; Décret n° 68-DF-478 & Acte
de naissance, preuve de filiation & OK \\
7 & \textbf{Contrat de travail (CDI/CDD)} & Employeur / MINTSS & Loi n°
92-007 du 14 août 1992 Code du travail (art. 23, 25) & Offert d'emploi,
bulletin de paie, etc. & OK \\
8 & \textbf{Immatriculation fiscale / NIU / attestation
d'immatriculation} & DGI (impots.cm, teledeclaration-dgi.cm) & Loi de
finances 2020 (art. L.1 bis) --- NIU obligatoire & Identité, activité,
adresse & OK \\
9 & \textbf{Carte grise / immatriculation véhicule} & Ministère des
Transports (SID) & Décret n° 79/341 du 03 sept. 1979 ; Décret n° 2020/PM
du 01 juin 2020 & Carte d'identité, facture d'achat, certificat de
dédouanement (Portrans.cm) & OK \\
10 & \textbf{Carnet de famille} & Officier d'état civil & Ordonnance n°
81/02 & Acte de mariage, actes de naissance des enfants & OK \\
11 & \textbf{Jugement supplétif d'acte de naissance} & Tribunal &
Ordonnance n° 81/02 (art. 23, 24, 33) & Preuve de la naissance tardive &
OK \\
12 & \textbf{Carte de séjour / autorisation de séjour} & DGSN & Décret
n° 2016/373 du 04 août 2016 (loi 97/012) & Passeport, justificatif de
séjour, moyens d'existence & OK \\
13 & \textbf{Légalisation de documents} & Ministère des Affaires
Étrangères (MAE) & Circulaire / pratique diplomatique & Document à
légaliser, identité & ATTENTION Texte peu accessible \\
14 & \textbf{Certificat de résidence / attestation d'hébergement} &
Mairie / chef de quartier & Règlement local & Preuve de domicile &
ATTENTION Pas de texte national clair ouvert \\
15 & \textbf{Diplôme / équivalence} & Ministère de l'Éducation & Texte à
identifier & Certificat, relevés de notes & ATTENTION Texte à
identifier \\
\end{longtable}

\textbf{Gap global} : les démarches ATTENTION (légalisation, attestation
de résidence, équivalence) illustrent l'\textbf{absence de textes
ouverts accessibles} pour certaines procédures → \textbf{à documenter
comme gaps} (valué, pas pénalisé).

\begin{center}\rule{0.5\linewidth}{0.5pt}\end{center}

\subsection{7. Plan d'action \& calendrier (urgent --- jusqu'au 10
novembre)}\label{plan-daction-calendrier-urgent-jusquau-10-novembre}

\textbf{Profil} : équipe technique, part de zéro, compétences de dev,
calendrier urgent.

\subsubsection{Phase 1 --- Fondations (J0--J3, \textasciitilde26--29
oct.)}\label{phase-1-fondations-j0j3-2629-oct.}

\begin{itemize}
\tightlist
\item[$\square$]
  Finaliser le concept et l'architecture (2 agents).
\item[$\square$]
  \textbf{Constituer la base de connaissances} (textes réels du §6,
  découpés en chunks JSON).
\item[$\square$]
  Squelette du code : Next.js (React) + FastAPI + Docker Compose.
\item[$\square$]
  Décision cloud vs local (démo cloud avec anonymisation ; version
  conforme local).
\end{itemize}

\subsubsection{Phase 2 --- MVP (J4--J7, \textasciitilde30 oct.--2
nov.)}\label{phase-2-mvp-j4j7-30-oct.2-nov.}

\begin{itemize}
\tightlist
\item[$\square$]
  \textbf{GOVIA-Assistant} : RAG + extraction d'intention + génération +
  pré-remplissage.
\item[$\square$]
  \textbf{GOVIA-Auditeur} : vérification des réponses vs textes, journal
  d'audit.
\item[$\square$]
  OCR des pièces (Tesseract).
\item[$\square$]
  Consentement explicite du citoyen.
\end{itemize}

\subsubsection{Phase 3 --- Livrables (J8--J10, \textasciitilde3--10
nov.)}\label{phase-3-livrables-j8j10-310-nov.}

\begin{itemize}
\tightlist
\item[$\square$]
  \textbf{Démo} (vidéo ≤ 7 min : version A vs B).
\item[$\square$]
  \textbf{Rapport technique} (5 pages max, template ITU).
\item[$\square$]
  \textbf{Dépôt GitHub} (code + base de connaissances ouverte).
\item[$\square$]
  \textbf{Mapping Y.3172} détaillé + \textbf{analyse d'écarts}.
\end{itemize}

\subsubsection{Risques \& mitigation}\label{risques-mitigation}

\begin{longtable}[]{@{}
  >{\raggedright\arraybackslash}p{(\columnwidth - 2\tabcolsep) * \real{0.5000}}
  >{\raggedright\arraybackslash}p{(\columnwidth - 2\tabcolsep) * \real{0.5000}}@{}}
\toprule\noalign{}
\begin{minipage}[b]{\linewidth}\raggedright
Risque
\end{minipage} & \begin{minipage}[b]{\linewidth}\raggedright
Mitigation
\end{minipage} \\
\midrule\noalign{}
\endhead
\bottomrule\noalign{}
\endlastfoot
Temps trop court & Prioriser le MVP (2 démarches : CNI + passeport) ;
compléter après. \\
API cloud indisponible / coût & Modèle local (Ollama) dès le départ. \\
Qualité FR/EN & Tests d'étalonnage ; écrémer les modèles. \\
Base de connaissances incomplète & Documenter les gaps ; ne pas
inventer. \\
\end{longtable}

\begin{center}\rule{0.5\linewidth}{0.5pt}\end{center}

\subsection{8. Livrables \& conformité au hackathon ---
checklist}\label{livrables-conformituxe9-au-hackathon-checklist}

\begin{longtable}[]{@{}
  >{\raggedright\arraybackslash}p{(\columnwidth - 2\tabcolsep) * \real{0.5000}}
  >{\raggedright\arraybackslash}p{(\columnwidth - 2\tabcolsep) * \real{0.5000}}@{}}
\toprule\noalign{}
\begin{minipage}[b]{\linewidth}\raggedright
Livrable / règle
\end{minipage} & \begin{minipage}[b]{\linewidth}\raggedright
Statut
\end{minipage} \\
\midrule\noalign{}
\endhead
\bottomrule\noalign{}
\endlastfoot
Partie 1 --- Solution technique (système fonctionnel : agent + RAG +
auditeur + démo) & OK Prévue \\
Partie 2 --- Ancrage réglementaire (dossier : mapping Y.3172 + analyse
d'écarts) & OK Prévue (§4.1) \\
\textbf{Date limite : les deux parties soumises avant le 10 novembre
2026} & ⚠️ Plan urgent (§7) \\
Mapping ITU-T Y.3172 détaillé & OK (§5.2 + annexe) \\
Base de connaissances RAG ouverte \& machine-readable & OK (§5.4) \\
Analyse d'écarts (gap analysis) & OK (§4.3, §6) \\
Comparaison version A vs B & OK (§5.5) \\
Politiques réelles et ouvertes (pas inventées) & OK \\
Données de cas d'usage synthétiques & OK \\
Rapport technique ≤ 5 pages & OK (à produire) \\
Vidéo démo ≤ 7 min & OK (à produire) \\
Dépôt GitHub & OK (à produire) \\
Aucune information confidentielle & OK \\
Équipe ≤ 4 personnes & ATTENTION À confirmer \\
Date limite 10 novembre 2026 & ATTENTION Plan urgent \\
\end{longtable}

\begin{center}\rule{0.5\linewidth}{0.5pt}\end{center}

\subsection{9. Recommandations \&
conclusion}\label{recommandations-conclusion}

GOVIA Cameroun est \textbf{faisable}, à condition de : 1.
\textbf{Respecter la Loi n° 2024/017} : consentement explicite, pas de
transfert transfrontalier non autorisé, pas de données biométriques →
privilégier un \textbf{modèle local}. 2. \textbf{Ne pas inventer de
politiques} : documenter chaque lacune (gap) --- c'est la \textbf{partie
la plus valorisée} du hackathon. 3. \textbf{Construire une base
multi-démarches} (CNI = exemple, pas la seule) avec des \textbf{textes
réels et ouverts}. 4. \textbf{Détailler le mapping Y.3172} et faire la
\textbf{démo version A vs B}. 5. \textbf{Travailler en urgence} sur un
MVP ciblé (CNI + passeport), livré avant le \textbf{10 novembre}.

\textbf{Prochaines étapes immédiates :} lancer la constitution de la
base de connaissances (chunks JSON des textes du §6) ; monter le
squelette du code ; décider du modèle (cloud pour démo avec
anonymisation, local pour conformité).

\begin{center}\rule{0.5\linewidth}{0.5pt}\end{center}

\subsection{Annexe A --- Mapping Y.3172 complété (tableau
détaillé)}\label{annexe-a-mapping-y.3172-compluxe9tuxe9-tableau-duxe9tailluxe9}

\begin{longtable}[]{@{}
  >{\raggedright\arraybackslash}p{(\columnwidth - 4\tabcolsep) * \real{0.3333}}
  >{\raggedright\arraybackslash}p{(\columnwidth - 4\tabcolsep) * \real{0.3333}}
  >{\raggedright\arraybackslash}p{(\columnwidth - 4\tabcolsep) * \real{0.3333}}@{}}
\toprule\noalign{}
\begin{minipage}[b]{\linewidth}\raggedright
Nœud
\end{minipage} & \begin{minipage}[b]{\linewidth}\raggedright
Description détaillée GOVIA
\end{minipage} & \begin{minipage}[b]{\linewidth}\raggedright
Politique applicable (exemples)
\end{minipage} \\
\midrule\noalign{}
\endhead
\bottomrule\noalign{}
\endlastfoot
SRC & Prompt citoyen + pièces (images) + base politiques (JSON) & Loi
2024/017 (finalité, minimisation) \\
C & Conversation + OCR (Tesseract) + indexation politiques & Loi
2024/017 (sécurité des données) \\
PP & Nettoyage, détection langue, \textbf{anonymisation}, normalisation
formats & Loi 2024/017 (art. 33 données sensibles, anonymisation) \\
M & LLM (Ollama/Groq) + RAG (ChromaDB) + extraction intent +
pré-remplissage & Décret 2016/375 (CNI, pièces) ; ordonnance 81/02 (état
civil) \\
P & Chaque étape cadrée par un texte cité & Loi 2024/017 (consentement,
DPIA) ; décret 2016/375 ; loi 92-007 \\
D & Réponse + formulaire pré-rempli + rapport d'audit vers le citoyen &
Loi 2010/012, 2010/021 (signature électronique, validation citoyenne) \\
SINK & Citoyen valide/soumet ; auditeur reçoit le rapport & Décret
2016/375 (validation finale) ; responsabilité du citoyen \\
\end{longtable}

\subsection{Annexe B --- Références des textes juridiques
(ouverts)}\label{annexe-b-ruxe9fuxe9rences-des-textes-juridiques-ouverts}

\begin{itemize}
\tightlist
\item
  Décret n° 2016/375 du 04 août 2016 (CNI) --- DGSN / prc.cm
\item
  Décret n° 2021/347 du 17 juin 2021 (Passeports) --- legecam.cm /
  prc.cm
\item
  Ordonnance n° 81/02 du 29 juin 1981 (État civil) ---
  droitcamerounais.info
\item
  Loi n° 2011/011 du 06 mai 2011 (délais état civil) --- camerlex.com
\item
  Loi n° 1968-LF-3 du 11 juin 1968 (Code de la nationalité) ---
  refworld.org / ILO
\item
  Loi n° 92-007 du 14 août 1992 (Code du travail) --- ILO /
  droitcamerounais.info
\item
  \textbf{Loi n° 2024/017 du 23 décembre 2024 (Protection des données)}
  --- Présidence / mathihis
\item
  Loi n° 2010/012 du 21 décembre 2010 (Cybersécurité) --- antic.cm
\item
  Loi n° 2010/021 du 21 décembre 2010 (Commerce électronique) ---
  mincommerce.gov.cm
\item
  Décret n° 79/341 du 03 septembre 1979 (Immatriculation véhicules) ---
  hacig.fr.gd
\item
  Décret n° 2020/PM du 01 juin 2020 (Cartes grises sécurisées) ---
  spm.gov.cm
\item
  Décret n° 2016/373 du 04 août 2016 (Étrangers) --- lc-doc.com
\item
  Procédure IDCAM --- idcam.cm ; Procédure PassCam --- passcam.cm
\item
  FGI 2026 / E-Gov'A 2026 --- contexte DPI
\end{itemize}

\subsection{Annexe C --- Scénario de démo (version A vs
B)}\label{annexe-c-scuxe9nario-de-duxe9mo-version-a-vs-b}

\textbf{Question du citoyen (ambiguë) :} « Je veux renouveler ma carte.
»

\begin{longtable}[]{@{}
  >{\raggedright\arraybackslash}p{(\columnwidth - 4\tabcolsep) * \real{0.3333}}
  >{\raggedright\arraybackslash}p{(\columnwidth - 4\tabcolsep) * \real{0.3333}}
  >{\raggedright\arraybackslash}p{(\columnwidth - 4\tabcolsep) * \real{0.3333}}@{}}
\toprule\noalign{}
\begin{minipage}[b]{\linewidth}\raggedright
\end{minipage} & \begin{minipage}[b]{\linewidth}\raggedright
Version A (sans gouvernance)
\end{minipage} & \begin{minipage}[b]{\linewidth}\raggedright
Version B (gouvernée)
\end{minipage} \\
\midrule\noalign{}
\endhead
\bottomrule\noalign{}
\endlastfoot
Réponse & Directe : « Renouvellez votre CNI en ligne sur idcam.cm. »
(risque : la carte peut être la carte grise ou la carte de séjour) &
Demande clarification ; identifie les démarches possibles (CNI, carte
grise, carte de séjour) ; \textbf{cite les textes} (décret 2016/375,
décret 79/341, décret 2016/373) ; l'\textbf{auditeur} signale
l'ambiguïté et demande précision. \\
Traçabilité & Aucune & Références aux textes + journal d'audit \\
Conformité & Risque d'erreur non contrôlé & Chaque conseil justifié par
un texte ; consentement requis \\
\end{longtable}

\begin{center}\rule{0.5\linewidth}{0.5pt}\end{center}

\emph{Document généré le 26 octobre 2026 --- Étude de faisabilité
technique et juridique GOVIA Cameroun. Respect des contraintes et règles
du hackathon ITU AI for Good Lab.}

\end{document}
